%% ============================================================
%%  Response to Reviewers
%%  Article: "A Biased Random-Key Genetic Algorithm with
%%            Rolling Horizon for the Meal Delivery Routing Problem"
%%  MIC 2026 Submission
%% ============================================================
\documentclass[12pt]{article}
\usepackage[margin=2.5cm]{geometry}
\usepackage{hyperref}
\usepackage{enumitem}
\usepackage{xcolor}
\definecolor{revcolor}{RGB}{0,70,140}
\newcommand{\rev}[1]{\textcolor{revcolor}{\textit{Reviewer: #1}}}
\newcommand{\resp}{\noindent\textbf{Response:}\ }
\newcommand{\change}{\noindent\textbf{Change in manuscript:}\ }

\begin{document}

\noindent
\textbf{Response to Reviewers}\\
\textit{A Biased Random-Key Genetic Algorithm with Rolling Horizon
for the Meal Delivery Routing Problem}\\[6pt]

We thank the reviewers for their careful reading and constructive
comments. Below we address each point in turn.
Manuscript changes are described immediately after each response
and are highlighted in the revised version.

\bigskip
\hrule
\bigskip

% -------------------------------------------------------------------
\noindent\textbf{Comment 1.}
\rev{A relevant modeling doubt concerns the treatment of order
coverage exclusively as an objective function rather than as a
constraint. In practical meal delivery systems, service-level
requirements often impose that all (or a very high percentage of)
orders must be fulfilled within the planning horizon. Please discuss
this choice in the manuscript.}

\resp
We agree that service-level agreements (SLAs) are common in practice
and that a coverage \emph{constraint} (e.g., $\rho \geq 0.95$) would
be a natural formulation. However, imposing such a constraint
directly on the rolling-horizon problem creates two fundamental
difficulties in a real-time setting.
First, whether a target coverage can be achieved depends on future
orders and courier availability that are unknown at dispatch time;
a constraint that is feasible at the end of the day may be locally
infeasible at every intermediate window, making the sub-problem
unsolvable and preventing any assignment from being dispatched.
Second, the real instances used in this study contain large,
high-volume days (up to 59,167 orders) in which even the best
deterministic and evolutionary methods achieve $\rho \approx
34$--$46\%$ — values that reflect structural limitations of the
available courier fleet rather than algorithmic shortcomings.
Imposing an SLA constraint above this ceiling would again render
the problem infeasible.
Treating coverage as an \emph{objective} avoids these pitfalls:
it always yields a feasible dispatch, and the algorithms
maximize it subject to the hard operational constraints
(capacity, time windows, eligibility) of Section~3.2.
This choice is consistent with the benchmark literature
(Reyes et al.\ 2018; Ulmer et al.\ 2021), where coverage or
delivery rate is reported as the primary outcome metric, not
a constraint.

\change
A paragraph discussing this modeling choice has been added at the
end of Section~3.3 (Performance Metric).

\bigskip

% -------------------------------------------------------------------
\noindent\textbf{Comment 2.}
\rev{Is the coverage rate considered as an objective function
computed for each time window or for the entire instance?}

\resp
The coverage rate $\rho$ plays two distinct roles in the paper,
and we recognise that the manuscript was not sufficiently clear
about this distinction.

\begin{itemize}
  \item \textbf{Per-window fitness (BRKGA only).}
  Inside the BRKGA decoder, $\rho$ is computed for the
  assignments produced in the current window $w$
  ($\rho^w = |\mathcal{A}^w| / |O^w|$) and is used as
  the primary component of the fitness functions $f_{\text{cov}}$
  and $f_{\text{trav}}$.
  This per-window value guides the evolutionary search within
  each 2-minute budget.

  \item \textbf{End-of-day metric (all methods, Table~2).}
  The $\rho$ values reported in Table~2 and discussed throughout
  Section~5 are computed \emph{over the entire planning day}:
  $\rho = |\bar{O}|/|O|$, where $\bar{O}$ is the set of
  all orders successfully delivered by the end of the 1,440-minute
  horizon. This is the primary performance metric for comparing
  methods.
\end{itemize}

\change
The end of Section~3.3 and the opening of Section~5 have been
clarified to distinguish the per-window fitness use from the
end-of-day reporting metric.

\bigskip

% -------------------------------------------------------------------
\noindent\textbf{Comment 3.}
\rev{Please explain better the structure of the chromosome.}

\resp
The chromosome is a real-valued vector of length
$n_{O^w} + n_{C^w}$ partitioned into two segments:
(i)~\emph{order keys} ($x_1, \ldots, x_{n_{O^w}}$): one
real value in $[0,1]$ per pending order; the decoder sorts
orders by these keys in descending order, so an order with
a higher key is attempted for assignment before orders with
lower keys.
(ii)~\emph{courier keys} ($x_{n_{O^w}+1}, \ldots,
x_{n_{O^w}+n_{C^w}}$): one real value per available courier;
the decoder tries couriers for each batch in descending key
order, so higher-key couriers are preferred.
The combined ranking induced by the two segments fully
determines the greedy assignment produced by the decoder
(Algorithm~2), without requiring any repair operator to
restore feasibility.

\change
Section~4.2.1 (Chromosome Encoding) has been expanded with a
concrete example tracing the decoding of a small chromosome to
illustrate how order and courier keys jointly determine the
assignment outcome.

\bigskip

% -------------------------------------------------------------------
\noindent\textbf{Comment 4.}
\rev{One of my main doubts concerns the definition of the
parameters (individual, probability, population size, delta,
etc). How are they selected?}

\resp
The parameters fall into two categories:

\textbf{Operational parameters} ($\Delta$, $K$, $D_{\max}$)
are dictated by the platform's operating rules rather than
algorithmic choices: $\Delta = 2$ min matches the platform's
dispatch cycle; $K = 3$ is the vehicle capacity limit; and
$D_{\max} = 10$ km is the platform's maximum courier-to-store
eligibility radius.

\textbf{BRKGA hyper-parameters} ($N_p$, $G_{\max}$, $N_e$,
$N_m$, $p_e$) were set by a two-step process.
First, canonical BRKGA defaults from
Gon\c{c}alves and Resende~(2011) were used as a starting
point: $p_e = 0.7$, $N_e / N_p \approx 0.15$--$0.20$,
$N_m / N_p \approx 0.10$--$0.15$.
Second, the values were tightened jointly so that the
worst-case per-window runtime on our largest instances
(up to 59,167 orders and 16,782 couriers) remains strictly
below $\Delta = 120$ seconds on the test hardware
(Intel Core i7, 16 GB RAM).
The binding constraint is the decoder cost, which is
$O(N_p \cdot G_{\max} \cdot |O^w| \cdot K)$ per window;
reducing $N_p$ or $G_{\max}$ below 10 produced noticeably
worse coverage on medium instances in pilot runs, while
increasing them caused occasional budget overruns on large
instances.
The final values ($N_p = 10$, $G_{\max} = 10$,
$N_e = 2$, $N_m = 2$, $p_e = 0.7$) represent the
Pareto-optimal trade-off between solution quality and
runtime safety observed in those pilot runs.

\change
A new paragraph ``Hyper-parameter calibration'' has been added
in Section~4.2, before the chromosome encoding, explaining the
rationale above.

\bigskip

% -------------------------------------------------------------------
\noindent\textbf{Comment 5.}
\rev{Why does the time remaining after the end of the genetic loop
remain unused? Why don't you also use a maximum number of
generations as a stop criterion?}

\resp
Two points of clarification:

First, $G_{\max} = 10$ \emph{is} the primary stopping criterion:
the genetic loop runs for exactly 10 generations and then stops,
regardless of how much wall-clock time has elapsed.
The 2-minute wall-clock budget ($\Delta$) is a \emph{safety
ceiling}, not the loop termination condition;
Algorithm~3 in the manuscript notes that the best solution
available at deadline $\Delta$ would be dispatched if the
loop were to overrun — but our parameterisation ensures
this never happens in practice.

Second, we deliberately do not reuse the remaining wall-clock
time for additional generations once $G_{\max}$ is reached.
The primary reason is \emph{predictability}: variable-length
runs would produce non-deterministic solution quality and
make the system harder to validate and maintain in production.
A fixed $G_{\max}$ also makes the per-window runtime
analytically bounded (independent of random convergence),
which is important for a real-time guarantee.
That said, adaptive stopping (e.g., early termination on
fitness stagnation, or additional generations when time
permits) is a natural direction for future work.

\change
Algorithm~3 and its surrounding text have been revised to
clearly state that $G_{\max}$ is the stopping criterion and
that the wall-clock budget serves only as a hard safety
ceiling. A sentence on why remaining time is not reused
has been added.

\bigskip

% -------------------------------------------------------------------
\noindent\textbf{Comment 6.}
\rev{In Table 2, the execution times in several cases are far
superior to the window size of 2 seconds. Please clarify this point.}

\resp
The execution times $T_{\text{BC}}$, $T_{\text{BT}}$, and
$T_{\text{BML}}$ reported in Table~2 are \emph{total daily
execution times} (in seconds), accumulated over all rolling-horizon
windows of a full 24-hour operating day — not per-window times.
A planning day of 1,440 minutes with $\Delta = 2$ minutes
contains up to 720 windows; summing the per-window runtimes
(each $\leq 120$ s) naturally yields totals in the range of
hundreds to thousands of seconds for large instances.
For example, Instance~0 (56,420 orders) has $T_{\text{BC}} =
3{,}143$ s, which corresponds to an average of $3{,}143 / 720
\approx 4.4$ seconds per window — well within the 120-second
per-window budget.
The 2-minute wall-clock constraint is therefore respected on
every window; Table~2 simply aggregates these times over
the day for efficiency comparison.

\change
The caption of Table~2 and the paragraph introducing it in
Section~5.1 have been updated to explicitly state that
$T$ values are \emph{total daily} execution times and to
note the per-window average for clarity.

\bigskip

% -------------------------------------------------------------------
\noindent\textbf{Comment 7.}
\rev{In the abstract, you state that our approach outperforms the
Business Model Lookahead and the Business Model by eight and fifteen
points; while at the end of the introduction, you state that the
percentage is ``eleven to thirteen''. Please check these sentences.}

\resp
Thank you for catching this inconsistency.
The figures are not contradictory — they refer to different
reference methods and instance scopes — but the introduction
text contained an error in the method name:

\begin{itemize}
  \item \textbf{Abstract}: ``eight and fifteen points'' refers to
  the \emph{overall} average advantage of the BRKGA variants over
  the Business Model Lookahead ($59.2\% - 51.4\% = 7.8 \approx 8$
  pp) and the plain Business Model ($59.2\% - 44.1\% = 15.1 \approx
  15$ pp), computed across all 24 instances.

  \item \textbf{Introduction}: ``eleven to thirteen percentage
  points'' correctly refers to the advantage over the
  \emph{Business Model Lookahead} on Small and Medium instances
  only (Small: $64.8\% - 53.6\% = 11.2$ pp;
  Medium: $59.1\% - 46.0\% = 13.1$ pp).
  The original text mistakenly wrote ``Business Model'' instead
  of ``Business Model Lookahead''.
\end{itemize}

\change
The last bullet of the contributions list in Section~1 has been
corrected to read ``Business Model Lookahead'' instead of
``Business Model'' for the eleven-to-thirteen percentage points
figure.

\bigskip

% -------------------------------------------------------------------
\noindent\textbf{Comment 8.}
\rev{Please clarify the definition of ghost kitchens used in the
literature review.}

\resp
Ghost kitchens (also called dark kitchens or virtual restaurants)
are food-preparation facilities that operate exclusively for
delivery orders, with no dine-in service and no fixed storefront
visible to the public.
Because their physical location can be chosen freely and updated
dynamically, they introduce an additional layer of spatial
flexibility to the MDRP: the platform must decide not only which
courier serves an order, but also from which kitchen that order
should be prepared and picked up.
The reference cited (Neria et al., \textit{Transportation Science}
2025) formalises this extension and shows that jointly optimising
kitchen and courier assignments improves coverage substantially.

\change
A brief parenthetical definition of ghost kitchens has been added
in Section~2.1 immediately after the citation of Neria et al.

\bigskip
\hrule
\bigskip

\noindent
We hope that the above responses and the corresponding manuscript
revisions satisfactorily address the reviewers' concerns.
We remain available for further clarification.

\end{document}
